\subsection{A regional observable that continues beyond first ignition}
On the two development woodpile edits these quarter crossings required 383 and 833 distinct local ignitions. The complete per-step integer count and mass curves matched across 180 steps, including 13 audited repair boundaries across both arms. Compared with the floating-point matte the integer version returned the same threshold frames. Times for complete ordinary, restart and repair were 275, 396 and 243 ms on A/B, and 588, 982 and 680 ms on B/A. Deferring 72 steps of material on the reverse edit made repair about 15.5\% slower than complete stepping, even as it saved the gas rerun relative to restart.

\subsection{Choosing between two strokes from the same shot}
A new seed and a fixed first brush define one common woodpile shot at frame 180. We registered a light second stroke with radius 0.06 and a full stroke with radius 0.12 around the same B center. They directly affected 74 and 288 sites respectively. The exterior selection and its initial 2,371 eligible sites were identical for both alternatives. The stipulated preference was to choose the lighter stroke if at least one quarter of target opacity mass ignited by 90 steps, otherwise the full stroke if it met the same deadline. The primary view and a second calibrated sensitivity view were fixed before the new trajectories.

In camera 0 the light stroke reached 11.69\% by the deadline. The full stroke crossed at $+44$ and reached 34.83\%. In camera 124 neither reached 0.08\%. All 48 replies and twelve 90-step curves matched the native reference. Both strokes encountered service boundaries. For the primary view, summing each candidate's median gives 739 ms for ordinary stepping, 922 ms for restart and 616 ms for repair. These estimates omit the outer checkpoint and restore between candidates. The following test measures the complete decision.
